\subsection{Reading Notes: Scattering processes and resonances from lattice QCD}
link: chrome-extension://efaidnbmnnnibpcajpcglclefindmkaj/https://arxiv.org/pdf/1706.06223\\\\
 another paper (or talk?) that goes into some more detail on this stuff:chrome-extension://efaidnbmnnnibpcajpcglclefindmkaj/https://iopscience.iop.org/article/10.1088/1742-6596/770/1/012024/pdf\\\\
Section 1: Introduction\\\\
- There are many exotic Hadrons which do not fit into the standard meson or baryon quark formations. These are unstable and appear as resonances, which can be inferred from bumps in scattering amplitudes.\\\\
- These resonances provide an opportunity to understand QCD more deeply in the non-perturbative low-energy region. \textbf{Why is the low energy region non-perturbative? does this have to do with the fine structure constant for gluon interaction?}\\\\
- Some motivation for understanding these resonances: "probes of the electroweak sector of the standard model consider processes which feature hadron resonances," basically these hadron resonances could improve our understanding of other regions of the standard model, and therefore have important application explaining actual physical processes. \\\\
- Lattice QCD is the only option available to understand these states, as it is the only non-perturbative approach to finding new states from first principles. \\\\
- The article discusses the use of lattice QCD to gain information about the properties of these exotic hadron resonances. \\\\
- an iconic resonance example is the $\sigma/f_0(500)$ meson which couples to the $\pi\pi$ scattering channel. It is unusual because the appearance of $\sigma$ was somehow unexpected, and it has a decay width larger than its mass, where decay width and mass are measured in units of energy. Decay width is related to decay time via $\Gamma =\hbar/ \tau$, where $\tau$ is the decay time \textbf{Decay width larger than mass? What?}\\\\
- Other examples of resonance tend to be smaller in width (resonance graphs have decay width vs mass).\\\\
- Lattice QCD considers quark and gluon fields on a discrete grid and samples many configurations of particles on the grid with probabilities determined by the Lagrangian of QCD. This has an associated error, which can be reduced by increasing the computing time. \\\\
- Most Lattice QCD calculations have dealt with stable hadrons rather than resonances, even getting the correct physical values with high accuracy (including small QED effects).\\\\
- Resonances appear as pole singularities in scattering amplitudes. \textbf{How? What does it mean to have a complex scattering amplitude? is it like a pole in the wavefunction?}.\\\\
- A current goal is to make a tower of excited hadron states for a single set of quantum numbers. LQCD can find up to a dozen of these states for a single quantum number channel.\\\\
Section 2: Resonances, Composite Particles, and Scattering Amplitudes\\\\
-Asymptotic states of QCD are the stable hadrons. \textbf{Why are they called asymptotic? In what way can they be described as asymptotes?}\\\\
- The paper investigates composite particles, which can be considered as a "dynamical enhancement in hadron scattering amplitudes" or as different constructions of quarks and gluons. Some of these are actually stable and are called bound states.\\\\
- An example of a bound state is a deuteron, which has the quantum numbers of a proton and neutron in an S-wave neutron ($l = 0$?), with fractional coupling to a D-wave neutron ($l = 2$?). Resonances of particles with longer lifetimes can be seen as "bump-like enhancements", or spikes,  in scattering amplitudes.\\\\
- scattering amplitude is given as ${\displaystyle f(\mathbf {k} _{f},\mathbf {k} _{i})=-{\frac {\mu }{2\pi \hbar ^{2}}}\int \psi _{f}^{*}V(\mathbf {r} )\psi _{i}d^{3}r}$ (from Wikipedia)\\\\
- It is related to cross section by ${\displaystyle d\sigma =|f(\theta )|^{2}\;d\Omega .}$ (Wikipedia)
- We consider a collision of two spinless particles, with energy and momentum $(E, \textbf{P})$ and center of momentum energy $E^\star$ .\\\\

